\section{Dominators for $\nu$ and the mean constant $m>1$}\label{sec:matching-mean}

\paragraph{What this section proves and how to read it.}
This section proves (CP6) and (CP7): it constructs positive laws
$M_{{\rm L},\Blk},M_{{\rm R},\Blk}$ dominating $\nu_\Blk$ at $\iota$
(called \emph{matching envelopes} in (CP6); they are unrelated to
the matching records $E_0,E_1$---in prose we call them
\emph{dominators for $\nu$}), with total price $T$ entering
$J=Be^T(R_I/2)^6$ through $e^T$, and it proves the mean energies
exceed $mY$ with $m>1$ (a lower bound \emph{added} to the energy,
hence $(1-m)/64<0$ in $\Theta$, unlike the subtracted savings $D$).
The \emph{side functions} below collect the diagonal weights; only
their final values are used, and the long list of intermediate
constants in \cref{def:side-function} that no estimate uses can be
skipped---the theorem statements cite only the assembled $T$ and
$m$. Reading-order note: $T$ includes upper-range prices proved in
\cref{sec:upper-comparison}; the per-block dominator construction
here is independent, the total $T$ is assembled in
\cref{sec:closing}. Inputs: \cref{sec:middle-frame}. Outputs:
(CP6),(CP7) with baseline $T<0.60760$, $m>1.58994$.
\begin{itemize}
\item For every atom $\Blk$ it constructs two positive measures
 $M_{{\rm L},\Blk}$ and $M_{{\rm R},\Blk}$, the \emph{matching envelopes}. The
 comparison law sits in a positive semidefinite $2\times2$ block
 $\bigl(\begin{smallmatrix}M_{{\rm L},\Blk}&\nu_\Blk\\
 \nu_\Blk^*&M_{{\rm R},\Blk}\end{smallmatrix}\bigr)\succeq_\iota0$ at the reserved
 cut, and $M_{s,\Blk}-\Nbr_\Blk$ is a positive reserved-diagonal measure. For each $s\in\{{\rm L},{\rm R}\}$, the sum over atoms $\sum_\Blk\operatorname{price}(M_{s,\Blk}-\Nbr_\Blk)$ is at most the matching exponent $T$ of the matching term
 $J=Be^T(R_I/2)^6$ (\cref{prop:record:matching}).
\item The proof of \cref{thm:comparison} uses the envelopes twice. Through the
 index-support lemma (\cref{lem:record:common-mean}) it compares the signed
 index form $J_\nu$ with $Y_\nu$, up to the excess of the random-inclusion
 energies $V_s$ of the envelopes over their mean energies $e_s$
 (\cref{eq:record:cm-selected}); through \cref{lem:envelope-energy} it bounds
 $V_s$ by $2^6J$. The mean energies are bounded below in
 \cref{subsec:record:positive_means}; this is the mean credit $m$ of $\Theta$
 (\cref{prop:record:mean-credit}, \eqref{eq:mean-credit}).
\item The closing (\cref{cert:closing}) evaluates $T$ and $m$.
\end{itemize}

Every record change of the comparison law except the two two-layer fold
records $F_1,F_2$ of a middle block already has such a bracket from its norm
envelopes (\cref{def:record:comparison-law}). For the matching records the
envelopes use chosen diagonals, which dominate the norm envelopes and are paid
by the same prices.

For $F_1,F_2$ the norm envelopes would enter $T$ with the full
fold norm price, which is too large for $\Theta<1$, and the
bracket is built
directly (\cref{subsec:record:direct_gram}). The reference record and its
signed comparison record are bounded jointly by a positive $2\times2$ Gram at
$\iota$, whose diagonal shifts are multiples of the record's own one-layer
companion. The dominator $\Nbr$ already contains that companion
once, with coefficient $v_{\rm D}$ (\cref{lem:record:gram}). The joint Gram
needs a diagonal of coefficient about $2.39\,w_{\rm D}$ on each side (the signed
coefficient $\bar H_Xd_\tlaw$), and only the excess
$k_s=w_{\rm D}d_s-v_{\rm D}$ is priced. The copy of the companion already present in $\Nbr$ is not counted again in that diagonal.

Each chosen side is an auxiliary side, given by a certified formula, plus a nonnegative rounded increment, and is at least the signed Gram coefficient $\bar H_Xd_\tlaw$. So only nonnegative increments are added to the
auxiliary price (\cref{subsec:signed-sides}). The envelopes and $T$ are
assembled in \cref{subsec:matching-envelopes}, the common mean of the
selected law in \cref{subsec:cm-selected}, and the mean credit in
\cref{subsec:record:positive_means}.

\subsection{A joint $2\times2$ block for reference and comparison records}\label{subsec:record:direct_gram}

The matching bracket uses the reference and the signed comparison record
jointly. Its leading part is a small matrix in the four actual features
$V_0,A_1/r,A_2/r,A_3/r$; it does not require a second calculation of their
individual spectra. We first isolate that matrix calculation, then state the
arithmetic side function used to price the resulting diagonals.

\begin{lemma}[The leading Gram]\label{lem:record:direct-gram}
Let $\mathcal H$ be a Hilbert space and let $V,A_1,A_2,A_3\colon H_0\to\mathcal H$
satisfy $\|(V,A_1/r,A_2/r,A_3/r)\|^2\le\Gamma$, where $r>0$.
For $\xi\in\R^3$ put $A_\xi=\sum_i\xi_iA_i$ and $\mathbf e=(1,0,0)$.
For $u,v\in\R^3$ define
\[
 L_{\rm old}=-(A_{\mathbf e}V^*+VA_{\mathbf e}^*),\qquad
 L_{u,v}=A_uV^*+VA_v^*,\qquad
 G(u,v)=\begin{pmatrix}L_{\rm old}&L_{u,v}\\L_{u,v}^*&L_{\rm old}\end{pmatrix}.
\]
For $\tside>0$ let
\[
 \Upsilon_{\tside}(u,v)=
 \begin{pmatrix}-\mathbf e/\tside&u\\v&-\tside\mathbf e\end{pmatrix},\qquad
 s(\tside;u,v)=\|\Upsilon_{\tside}(u,v)\|.
\]
This is a real $6\times2$ matrix, and
\begin{gather*}
 s(\tside;u,v)^2=\lambda_{\max}
 \begin{pmatrix}\tside^{-2}+\|v\|^2&-u_1/\tside-\tside v_1\\
 -u_1/\tside-\tside v_1&\tside^2+\|u\|^2\end{pmatrix},\\
 G(u,v)\succeq-r\Gamma s(\tside;u,v)
 \operatorname{diag}(\tside I,\tside^{-1}I).
\end{gather*}
For the signed arm vector $\mathbf a_\tlaw=(\tlaw,1,1-\tlaw)$, put
$u_\tlaw=\mathbf a_\tlaw-\mathbf e$ and
$v_\tlaw=-\mathbf a_\tlaw-\mathbf e$. Then
\[
 L_{u_\tlaw,v_\tlaw}=L_{\rm old}
 +A_{\mathbf a_\tlaw}V^*-VA_{\mathbf a_\tlaw}^*,\qquad
 G(u_\tlaw,v_\tlaw)\succeq-r\Gamma s_\tlaw I,
\]
where
\[
 s_\tlaw=\bigl(4-2\tlaw+2\tlaw^2+2\sqrt{1+\tlaw^2}\bigr)^{1/2},
 \qquad s_0=\sqrt6.
\]
The statement holds directly in the actual half-marginal space; the adjoints
and the frame bound are always those of $\mathcal H$.
\end{lemma}

\begin{proof}
Let $\Sigma=(V,A_1/r,A_2/r,A_3/r)$ and
$\mathbf S=\operatorname{diag}(\Sigma,\Sigma)$. In each copy of $H_0^4$
separate the centre coordinate $V$ from the three arm coordinates. After
conjugating $G(u,v)$ by
$\Delta_{\tside}^{-1/2}$, where
$\Delta_{\tside}=\operatorname{diag}(\tside I,\tside^{-1}I)$, its coefficient
matrix in these eight features, ordered with both centres first, is
\[
 \mathbf M=\begin{pmatrix}0&\Upsilon_{\tside}^{\mathsf T}\\
 \Upsilon_{\tside}&0\end{pmatrix}.
\]
Indeed the old diagonal stars acquire the coefficients
$-\mathbf e/\tside$ and $-\tside\mathbf e$, while the off-diagonal arms retain
$u,v$. Thus the conjugated operator is
$r\mathbf S(\mathbf M\otimes I)\mathbf S^*$, with the same coordinate
permutation in $\mathbf S$. The least eigenvalue of $\mathbf M$ is
$-\|\Upsilon_{\tside}\|$, so it is at least
$-r s(\tside;u,v)\mathbf S\mathbf S^*\succeq-r\Gamma s(\tside;u,v)I$.
Conjugating back proves the order. The displayed $2\times2$ matrix is
$\Upsilon_{\tside}^{\mathsf T}\Upsilon_{\tside}$.

At the signed point, linearity of $A_\xi$ gives the identity for
$L_{u_\tlaw,v_\tlaw}$, and at $\tside=1$ the last matrix is
\[
 \begin{pmatrix}4+2\tlaw^2&2\\2&4-4\tlaw+2\tlaw^2\end{pmatrix}.
\]
Its largest eigenvalue is $4-2\tlaw+2\tlaw^2+2\sqrt{1+\tlaw^2}$.
\end{proof}

In the application $\mathcal H=L^2(m_0)$ and the features are those of
\cref{lem:actual-features}, with $V=V_0$ and $\Gamma=\Gamma_A$.
The reference contributes $L_{\rm old}$ on the diagonal; the signed law
contributes its old leading part plus its skew leading difference off the
diagonal (\cref{lem:old-f-law}(f), \cref{prop:signed-bounds}(c)). This is why the
Gram gives a joint bracket and not only a norm bound for the change. At
$\tlaw=127/500$, $s_\tlaw=2.38422727\ldots$, compared with $\sqrt6=2.4494897\ldots$
at the mirror point. Adding the two remainders gives the actual coefficient
\[
 d_\tlaw=s_\tlaw\Gamma_A+(e_o+e_{\rm sg})r,
\]
which is below $2.39033202$ at $Q_3$. The block argument is
\cref{prop:signed-sides} below.

\paragraph{Auxiliary sides and the certified price.}
The Gram in the actual space $L^2(m_0)$ is supplied by the signed law. The numerical price is stored
as an auxiliary side plus a nonnegative increment to that Gram coefficient.
We retain the complete arithmetic definition of the auxiliary side because it
fixes this split and its outward roundings. Its parameters are inherited from
a probability mixture of four sign patterns, but no law of the construction
uses that mixture; no operator estimate for it is needed. In particular the
coefficients below are definitions of numbers, not assertions about another
comparison law.

\begin{definition}[The side function]\label{def:side-function}
For a prime $p\equiv3\pmod4$ with $q=1/p\le2^{-20}$, put
$r=p^{-1/2}$, $Z=(1-q)(1-2q)$,
$h_0=1-14q$ and $e_0=2q/((1-3q)(1-4q))$, as in
\cref{subsec:middle-records:notation}. Fix
\[
 (w_{++},w_{+-},w_{-+},w_{--})=(5803723,2409463,1786814,0)/10^7,
\]
\[
 (\bar b,\bar d,\bar f)=(b_0,d_0,f_0)=(.6426372,.5181074,.1607446),\qquad
 \tside=31/32,
\]
where $\bar b=\sum_{(\beta,\delta)\in\{\pm1\}^2}\beta
w_{\beta\delta}$,
$\bar d=\sum_{(\beta,\delta)\in\{\pm1\}^2}\delta
w_{\beta\delta}$ and
$\bar f=\sum_{(\beta,\delta)\in\{\pm1\}^2}\beta\delta
w_{\beta\delta}$.
Set $u=(-\bar b,1,\bar b)$, $v=(-1,-\bar d,-\bar d)$,
$s(\tside)=s(\tside;u,v)$ as in \cref{lem:record:direct-gram}, and
\begin{gather*}
 \gamma=\left(1+\frac{q(1+2r)}{1-2q}\right)^{1/2},\qquad
 c_{\rm lead}=\gamma(\|u\|_2+\|v\|_2),\\
 F_{\rm frame}=\frac{2+|\bar b+1|}{1-q}
 +\frac{|1-\bar d|q}{(1-q)(1-2q)}
 +\frac{2|\bar b-\bar d|q^2}{(1-q)(1-2q)}
 +\frac{|2-\bar b+\bar d|r}{1-q},\\
 \zeta_{++}=\frac{6q^2}{(1-4q)(1-5q)},\qquad
 \zeta_{+-}=\frac{4q}{(1-4q)(1-5q)},\\
 \zeta_{-+}=\frac{4q}{(1-4q)(1-5q)(1-7q)},\qquad
 \zeta_{--}=\frac{8q}{(1-4q)(1-5q)(1-7q)},\\
 C_{\rm old}=\frac{19}{Zh_0},\qquad
 C_{\rm ns}=2|\bar b|+1+2|\bar d|+4|\bar f|,\qquad
 n_*=\frac{1+8r+(18/Z+144)q}{h_0},\\
 C_z=\frac{n_*}q\sum w_{\beta\delta}\zeta_{\beta\delta},\qquad
 C_{\rm new}=\frac{2C_{\rm ns}/Z+144+|\bar b-\bar d|}{h_0}+C_z,\\
 k_{\rm raw}=\frac{rc_{\rm lead}}{h_0}+qC_{\rm new},\qquad
 \eta_1=\frac{(4w_{++}+6w_{+-}+4w_{-+}+8w_{--})q}{h_0},\\
 \rho_1=\frac{k_{\rm raw}+\eta_1^2}{1-\eta_1},\qquad
 \eta_\star=\frac{e_0+\eta_1}{1-\eta_1},\qquad
 \phi_2=\frac{\eta_\star}{1-\rho_1},\qquad
 K_\infty=\frac4{h_0^{1/2}(1-\eta_1)},\\
 \phi_\infty=\eta_\star+K_\infty\phi_2,\qquad
 C_{\rm rep}=\frac{2\phi_\infty}q\Bigl(k_{\rm raw}+\frac{14q}{h_0}\Bigr),
 \qquad C'=C_{\rm new}+C_{\rm rep}.
\end{gather*}
With these numbers define
\begin{gather*}
 c_{\rm com}=\frac{rF_{\rm frame}}{h_0}
 +r\Bigl(C_{\rm old}+C'+\frac{2e_0^2}q\Bigr),\\
 d^m_{\rm L}=\frac{\gamma s(\tside)\tside/h_0+c_{\rm com}}{1-e_0},\qquad
 d^m_{\rm R}=\frac{\gamma s(\tside)\tside^{-1}/h_0+c_{\rm com}}{1-e_0},\\
 k_0=\frac{2\gamma/h_0+r(C_{\rm old}+14/h_0+2e_0/q+e_0^2/q)}{1-e_0},\\
 k_1=\frac{c_{\rm lead}/h_0+r(C'+14/h_0+2e_0/q+e_0^2/q)}{1-e_0},\\
 \bar k=\max(k_0,k_1).
\end{gather*}
For a middle piece starting at $X$, evaluate at $q_X=1/X$ and
$r_X=1/\lfloor\sqrt X\rfloor\ge X^{-1/2}$, and put
$n_X=\lfloor L_{\rm loc}/\log_-X\rfloor$, $L_{\rm loc}=585$
(\cref{lem:record:local-domain}). The \emph{auxiliary sides} are
\[
 d^{\rm aux}_s(X)=e^{n_Xe_0(X)}
 \max\bigl(d^m_s(X),\,2\bar k(X)^2r_X\bigr),\qquad s\in\{{\rm L},{\rm R}\}.
\]
Every coefficient is evaluated as follows: the displayed
expression, exactly as written as a function of the independent
variables $(q,r)$ in the nondecreasing form of
\cref{lem:majorant} (each remainder $qc$ written as $r(rc)$), is
evaluated at $(q_X,r_X)$; no coefficient is used in a different
$(q,r)$ form. In particular a remainder evaluated at $q_X$ is
never divided by the larger value $r_X$.
The $33$ values and their gates, including $\eta_1<1$, $\rho_1<1$ and
$2\phi_\infty<1$, are those of
\cref{tab:data:sides,cert:record:matching}.
\end{definition}

At this auxiliary point $s(31/32)=2.0778741\ldots$,
$s\tside=2.0129405\ldots$ and $s/\tside=2.1449023\ldots$.
The chosen physical sides below are at least the signed Gram coefficient,
regardless of these smaller auxiliary values.

For the block argument, recall that $\nu_F$ is the uniform law of a half
$(x_0,x_1,x_2)$ with distinct labels and $x_0\to x_2$. The reference
record has density $K_0=\Psi_0/z_0$ relative to $\nu_F\otimes\nu_F$,
where $\Psi_0=I_\times(1+\mathsf A)(1+\mathsf B)$ and both half marginals
are $m_0=g_0\nu_F$. Here $h_0\le z_0\le1$ and $|g_0-1|\le e_0$
(\cref{lem:old-f-law}(a)); all Markov adjoints in the next lemma are in
$L^2(m_0)$.

\begin{lemma}[Unequal sides]\label{lem:record:asymmetric-gram}
Let $\Blk$ be a middle block of a piece starting at $X$, with least prime $b$,
$r_b=b^{-1/2}$ and at most $n_X=\lfloor L_{\rm loc}/\log_-X\rfloor$ primes. Let
$\Pi'_{F_1,\Blk}$ be a record placed by $F_1$ whose kernel is a finite real
combination
\[
 K'_\Blk=\sum_{c\in I}\alpha_c\bigotimes_{p\in\Blk}K'_{c,p},\qquad
 \sum_{c\in I}\alpha_c=1,
\]
of products of probability couplings $K'_{c,p}$ with both half marginals $m_0$
(a single product is the case $|I|=1$); write $T'_{c,p}$ for their Markov
operators on $L^2(m_0)$ and $T'_p=\sum_c\alpha_cT'_{c,p}$.
\begin{enumerate}[label=(\alph*)]
\item \emph{Enlargement.} If the block cut matrix of (b) is $\succeq_\iota0$ with
 sides $d_{\rm L},d_{\rm R}$, it is so with any sides $d'_{\rm L}\ge d_{\rm L}$,
 $d'_{\rm R}\ge d_{\rm R}$.
\item \emph{Block extension.} Suppose that, for nonnegative functions $c_{\rm L},c_{\rm R}$
 of $(q,r)$ that are nondecreasing in $(q,r)$ and a number $c_{\ge2}(X)\ge0$,
 \begin{itemize}
 \item at every prime $p\in\Blk$, on $L^2(m_0)^2$,
  $\bigl(\begin{smallmatrix}T_{\rm old}+r_pc_{\rm L}I&T'_p\\
  T'^*_p&T_{\rm old}+r_pc_{\rm R}I\end{smallmatrix}\bigr)\succeq0$, with
  $c_s=c_s(1/p,p^{-1/2})$ and $T_{\rm old}$ the Markov operator of $K_0$;
 \item on every exact support $S\subseteq\Blk$ of at least two primes, the
  restrictions to the sector $\mathcal H_S$ of the old block operator
  $\bigotimes_pT_{{\rm old},p}$ and of the new one $\sum_c\alpha_c\bigotimes_pT'_{c,p}$
  have norms summing to at most $c_{\ge2}(X)r_Xr_b$.
 \end{itemize}
 Then, as cut matrices at $\iota$ of measures on the labels of $\Blk$,
 \[
  \begin{pmatrix}\Pi_{F_1,\Blk}+r_bd_{\rm L}D^1_\Blk&\Pi'_{F_1,\Blk}\\
  \Pi'^{\,*}_{F_1,\Blk}&\Pi_{F_1,\Blk}+r_bd_{\rm R}D^1_\Blk
  \end{pmatrix}\succeq_\iota0
  \qquad\text{for}\quad
  d_s\ge H_\Blk\max\bigl(c_s(X),c_{\ge2}(X)r_X\bigr),
 \]
 where $H_\Blk$ is any bound for the density of $m_{0,\Blk}=\bigotimes_pm_{0,p}$
 relative to $\nu_F^{\otimes\Blk}$, for instance $e^{n_Xe_0(X)}$, or
 $\bar H_X=(1-n_Xe_0(X))^{-1}$ if $n_Xe_0(X)<1$; and $D^1_\Blk$ is the diagonal
 copy at $\iota$ of the product law $\nu_F^{\otimes\Blk}$ placed by $F_1$, the
 normalized one-layer companion of $F_1$ in $\widetilde\mu_\Blk$.
\item Let $\varsigma(x)=(-x_2,-x_1,-x_0)$ and $\Phi(x,y)=(\varsigma
 y,\varsigma x)$. Then $\varsigma$ preserves $\nu_F$ and $g_0$, $\Phi$
 preserves $K_0$, and $\Phi$ maps the edges $x_0\to y_2$,
 $y_0\to x_2$, $x_1\to y_0$, $x_0\to y_1$ to $x_0\to y_2$, $y_0\to x_2$,
 $x_2\to y_1$, $x_1\to y_2$. If the comparison record of $F_2$ is the
 reflected transpose of that of $F_1$, $K'_{2,c,p}(x,y)=K'_{c,p}(\varsigma
 y,\varsigma x)$ for every component and every prime, then (b) holds for
 $F_2$ with $d_{\rm L}$ and $d_{\rm R}$ interchanged.
\end{enumerate}
Consequently, if the \emph{side gates} $k_s:=w_{\rm D}d_s-v_{\rm D}\ge0$,
$s\in\{{\rm L},{\rm R}\}$, hold in the piece, the two fold records of a block
cost $k_{\rm L}+k_{\rm R}$ per unit $r_b^2$ on each side.
\end{lemma}

\begin{proof}
(a) The difference of the two cut matrices is the diagonal
$\operatorname{diag}((d'_{\rm L}-d_{\rm L})r_bD^1_\Blk,(d'_{\rm R}-d_{\rm R})r_bD^1_\Blk)$,
a positive reserved-diagonal measure, which is $\succeq_\iota0$ (\ref{it:record:E3}).

(b) The old block operator is $\bigotimes_pT_{{\rm old},p}$ and the new one is
the combination $T'_\Blk=\sum_c\alpha_c\bigotimes_pT'_{c,p}$; every prime factor
fixes the constants of $m_0$ on both sides, because all these couplings have
the actual marginal $m_0$ on both sides. By \cref{lem:sector}(a) and
\cref{rem:signed-mixture-sectors} both preserve the orthogonal sectors
$\mathcal H_S$ on which the factors in a set $S$ of primes are mean-zero and the
others constant; the old one is $\bigotimes_{p\in S}(T_{{\rm old},p}-\Jproj)$ there,
and the new one is $T'_p-\Jproj$ at the one coordinate $p$ if $S=\{p\}$. So the
block Gram is the orthogonal sum of its restrictions to the sectors
(\cref{lem:sector}(a)), and it suffices to treat each sector with the shifts
$\sigma_s=r_b\max(c_s(X),c_{\ge2}(X)r_X)$. On $\mathcal H_\emptyset$ it is
$\bigl(\begin{smallmatrix}1+\sigma_{\rm L}&1\\1&1+\sigma_{\rm R}
\end{smallmatrix}\bigr)\succeq0$, as $\sum_c\alpha_c=1$. On $S=\{p\}$ it is the
restriction of the prime Gram at $p$, which holds with the shifts
$r_pc_s(1/p,p^{-1/2})\le r_bc_s(X)\le\sigma_s$, by the monotonicity of $c_s$, as
$p\ge b\ge X$ and $r_X\ge X^{-1/2}$; larger shifts preserve it. For $|S|\ge2$ let
$O$ and $N$ be the restrictions of the old and new operators; $O$ is
self-adjoint, $\|O\|+\|N\|\le c_{\ge2}(X)r_Xr_b\le\sigma_s$, and
$\bigl(\begin{smallmatrix}O+\sigma_{\rm L}&N\\N^*&O+\sigma_{\rm R}\end{smallmatrix}\bigr)
\succeq\bigl(\begin{smallmatrix}(\sigma_{\rm L}-\|O\|)I&N\\N^*&(\sigma_{\rm R}-\|O\|)I
\end{smallmatrix}\bigr)\succeq0$, since $\sigma_s-\|O\|\ge\|N\|$. This gives the block
Gram on $L^2(m_{0,\Blk})^2$ with the shifts $\sigma_s$ times $I$.

In measure form, the cut forms at $\iota$ of $\Pi_{F_1,\Blk}$ and
$\Pi'_{F_1,\Blk}$ at real half-functions are the forms of $T_{{\rm old},\Blk}$ and
$T'_\Blk$ on $L^2(m_{0,\Blk})$, and that of the identity is the cut form of the
diagonal law $\operatorname{Diag}_\iota(m^{F_1}_{0,\Blk})$ of the actual half
marginal placed by $F_1$. Its density with respect to $D^1_\Blk$ is
$\prod_{p\in\Blk}g_0(x_p)\le H_\Blk$; for the two examples, $g_0\le1+e_0(X)$ at
every prime of the block (\cref{lem:old-f-law}(a)), and the block has at most
$n_X$ primes, so the density is at most $(1+e_0(X))^{n_X}\le e^{n_Xe_0(X)}\le
\bar H_X$. Replacing $\operatorname{Diag}_\iota(m^{F_1}_{0,\Blk})$ by
$H_\Blk D^1_\Blk$ therefore adds a positive reserved-diagonal measure, which is
$\succeq_\iota0$ (\ref{it:record:E3}).

(c) $\varsigma$ is an involution preserving $\nu_F$
(\cref{subsec:middle-records:notation}) and $g_0$ (\cref{lem:old-f-law}(a)).
Writing
$(x',y')=\Phi(x,y)$, a relation $s\to t$ between labels becomes $-t\to-s$;
this gives the edge table, and in particular $x_0\to y_2$ and $y_0\to x_2$
are fixed, so $K_0\circ\Phi=K_0$. The added edges of $F_1$
go to those of $F_2$ (\cref{cert:record:middle-geometry}). With
$(\mathcal U\varphi)(x)=\varphi(\varsigma x)$, unitary on $L^2(\nu_F)$ and on
$L^2(m_0)$, the hypothesis reads $T'_{2,c,p}=\mathcal UT'^{\,*}_{c,p}\mathcal U^*$
for every component, and $\mathcal UT_{\rm old}\mathcal U^*=T_{\rm old}$. Exchange
the two copies in the prime Gram of (b) and conjugate by
$\operatorname{diag}(\mathcal U,\mathcal U)$: this is the prime Gram of the $F_2$
record with $c_{\rm L}$ and $c_{\rm R}$ interchanged, the shifts being multiples of
the identity. On the block, $\bigotimes_p\mathcal U$ maps every sector
$\mathcal H_S$ onto itself, so the bound on the sectors of at least two primes
is carried over unchanged, and the density of $m_{0,\Blk}$ is $\varsigma$-invariant.
So the proof of (b) applies to $F_2$ with the sides interchanged.

The last assertion is the bookkeeping of the two records: by (b) and (c),
the layer-zero side carries $r_bd_{\rm L}$ for $F_1$ and $r_bd_{\rm R}$ for
$F_2$, the layer-one side the reverse, each record has weight $w_{\rm D}r_b$, and
the bracket already contains the one-layer companion of each record with the
weight $v_{\rm D}r_b^2$ (one fold copy, \cref{lem:record:gram}).
\end{proof}

\paragraph{The global substitution.}
For a middle block $\Blk$ let
$F^{(2)}_\Blk=w_{\rm D}r(\Pi_{F_1,\Blk}+\Pi_{F_2,\Blk})/\widetilde\mu_\Blk(1)$ be the reference's two
two-layer fold records, with their one-layer companions
$F_{{\rm fold},\Blk}=v_{\rm D}r^2(D^1_\Blk+D^2_\Blk)/\widetilde\mu_\Blk(1)$ (\cref{def:record:frame}(b),
where $D^j_\Blk$ is written $\operatorname{Diag}_\iota(\Pi^{(1)}_{F_j})$),
where here $\Pi_{F_j,\Blk}$ and $D^j_\Blk$ denote probability laws;
the common normalizing factor is displayed explicitly, unlike the
component notation of \cref{def:record:frame}(b). The law
$D^2_\Blk$ is defined as
$D^1_\Blk$ with $F_2$ in place of $F_1$; and put
\[
 \lambda_{{\rm res},\Blk}=\lambda_\Blk-F^{(2)}_\Blk,\qquad
 \Nbr_{{\rm res},\Blk}=\Nbr_\Blk-F^{(2)}_\Blk-F_{{\rm fold},\Blk}
 =\lambda_{{\rm res},\Blk}+H_\Blk .
\]
Outside the middle put $\lambda_{{\rm res},\Blk}=\lambda_\Blk$ and
$\Nbr_{{\rm res},\Blk}=\Nbr_\Blk$.

\subsection{The sides of the signed law}\label{subsec:signed-sides}

The joint $2\times2$ block of the preceding subsection needs, on
each of its two diagonal sides, a coefficient---the \emph{side}:
the multiple of the one-layer companion record placed on that
diagonal. For the signed comparison records of $F_1$ and $F_2$
the block is positive already with the two equal sides
$\bar H_Xd_\tlaw(X)$, where $d_\tlaw$ is the explicit coefficient
computed from the joint Gram matrix in
\cref{prop:signed-bounds}(d). The construction instead starts
from cheaper auxiliary sides given by a certified formula
(\cref{def:side-function}) and increases each by a nonnegative
increment to at least $\bar H_Xd_\tlaw(X)$; only the auxiliary
price and the increment are then charged. The proposition makes
this precise, piece by piece.

\begin{proposition}[The sides of the signed law]\label{prop:signed-sides}
For a middle piece starting at $X$ let $n_X=\lfloor L_{\rm loc}/\log_-X\rfloor$,
assume $n_Xe_0(X)<1$, put $\bar H_X=(1-n_Xe_0(X))^{-1}$, and let $d_\tlaw(X)$ and
$c_{\ge2}(X)$ be the constants of \eqref{eq:signed-constants} at $(q_X,r_X)$, with
$\tlaw=127/500$ (the functions of \cref{cert:signed-local} evaluate them at the
$33$ piece starts).
\begin{enumerate}[label=(\roman*)]
\item Let $\Blk$ be a middle block of the piece. The signed comparison record of
 $F_1$ (\cref{def:record:mixture}(c)) satisfies the hypotheses of
 \cref{lem:record:asymmetric-gram}(b) with $c_{\rm L}=c_{\rm R}=d_\tlaw$ and the
 support-two constant $c_{\ge2}(X)$, and $c_{\ge2}(X)r_X<d_\tlaw(X)$. Hence its block
 Gram holds with $d_{\rm L}=d_{\rm R}=\bar H_Xd_\tlaw(X)$, and so does that of $F_2$
 (\cref{lem:record:asymmetric-gram}(c)).
\item The \emph{increments}
 \begin{equation}\label{eq:signed-increment}
  \delta^+_{s,X}=\operatorname{up}\bigl([\bar H_Xd_\tlaw(X)-d^{\rm aux}_s(X)]_+\bigr),
  \qquad s\in\{{\rm L},{\rm R}\},
 \end{equation}
 where in the display $\bar H_Xd_\tlaw(X)$ is replaced by its certified upper endpoint $U_X$ from \cref{cert:signed-local}, and with
 $\operatorname{up}(x)=10^{-30}\lceil10^{30}x\rceil$, define the \emph{chosen sides}
 $d_s(X)=d^{\rm aux}_s(X)+\delta^+_{s,X}\ge\max(\bar H_Xd_\tlaw(X),d^{\rm aux}_s(X))$,
 for which the block Grams of (i) hold (\cref{lem:record:asymmetric-gram}(a)). In
 the piece $i$, starting at $X_i$, with the completion weights $w_{{\rm D},i}$,
 $v_{{\rm D},i}$ of \cref{lem:front12:reference}, the \emph{side coefficients}
 are
 \[
  k_{s,i}=w_{{\rm D},i}d_s(X_i)-v_{{\rm D},i}=k^{\rm aux}_{s,i}+w_{{\rm D},i}\delta^+_{s,i},
  \qquad k^{\rm aux}_{s,i}=w_{{\rm D},i}d^{\rm aux}_s(X_i)-v_{{\rm D},i},
 \]
 with $\delta^+_{s,i}=\delta^+_{s,X_i}$, and $k^{\rm aux}_{s,i}\ge0$ in every piece
 (\cref{cert:record:matching}); so the side gates hold. (These side coefficients
 are not the caps $k_{\rm L}$ of the layer cut L21 and $k_{\rm R}=k_{\rm E}$ of
 \cref{lem:front12:reference}(ii); in this section $k_{\rm L},k_{\rm R}$ always
 denote side coefficients.)
\item With the certified degree-two charges
 $\bar I_{2,i}=\operatorname{hi}(I_{2,i})\ge\sum_\Blk r_\Blk^2q_\Blk^{2a_0}$ of the pieces at
 the anchor $a_0$ (\cref{lem:record:block-charge}, \cref{def:closing-quantities}),
 \[
  \sum_i(k_{{\rm L},i}+k_{{\rm R},i})\bar I_{2,i}\le T'_K:=T^{\rm aux\prime}_K+T^{\rm inc\prime}_K,
 \]
 where $T^{\rm aux\prime}_K=\sum_i\operatorname{up}\bigl((k^{\rm aux}_{{\rm L},i}
 +k^{\rm aux}_{{\rm R},i})\bar I_{2,i}\bigr)$ and $T^{\rm inc\prime}_K=\sum_i
 \operatorname{up}\bigl(2w_{{\rm D},i}\operatorname{up}\bigl(\frac12(\delta^+_{{\rm L},i}
 +\delta^+_{{\rm R},i})\bigr)\bar I_{2,i}\bigr)$. Put $T_K=tT'_K$, with the repricing
 factor $t=e^{2(a-a_0)L}$, $L=570$ (\cref{def:closing-quantities}(1)).
\end{enumerate}
\end{proposition}

\begin{proof}
(i) \emph{The prime Gram.} Let $p\in\Blk$, so $p\ge X\ge Q_3>2^{24}$, and let
$T'_p=T[K_{\tlaw,p}]$ be the operator on $L^2(m_0)$ of the signed record at $p$
(the notation $T[\cdot]$ of \cref{prop:signed-bounds}), the combination
$P_{++,p}+\frac\tlaw2P_{-+,p}-\frac\tlaw2P_{--,p}$ of its three components
(\cref{def:record:mixture}(c)). These are probability couplings with both half
marginals $m_0$ (\cref{lem:signed-components}(d)), so
$T_{\rm old}$ and $T'_p$ fix the constants on both sides, and the decomposition
of $L^2(m_0)$ into the constants and the $m_0$-centred functions reduces the
prime Gram. On the constants it is
$\bigl(\begin{smallmatrix}1+dr&1\\1&1+dr\end{smallmatrix}\bigr)\succeq0$ for
$d\ge0$. On the centred functions, $T_{\rm old}=O_0+R_o$ with
$\|R_o\|\le e_oq$, where $O_0=-(A_{\mathbf e}V_{0,m}^*+V_0A_{\mathbf e}^{*_m})$ is the
natural leading part, with adjoints in $L^2(m_0)$ (\cref{lem:old-f-law}(f));
and $T'_p=O_0+L^m_{\tlaw,p}+R^{\rm new}_{\tlaw,p}$ with
$L^m_{\tlaw,p}=A_{\mathbf a_\tlaw}V_{0,m}^*-V_0A_{\mathbf a_\tlaw}^{*_m}$, the signed
leading difference with the unprojected pair feature $A_3$ and the adjoints of
$L^2(m_0)$, and $\|R^{\rm new}_{\tlaw,p}\|\le e_{\rm sg}q$
(\cref{prop:signed-bounds}(c)). The leading parts form the operator matrix
$G(u_\tlaw,v_\tlaw)$ of the synthesis form of \cref{lem:record:direct-gram}, with
$\mathcal H=L^2(m_0)$, $V=V_0$ and the features $A_i$ regarded as maps into
$L^2(m_0)$, whose synthesis has squared norm at most the frame $\Gamma_A$ of
\cref{lem:actual-features}(d). At $\tside=1$ it gives
$G(u_\tlaw,v_\tlaw)\succeq-r\Gamma_As_\tlaw I$. The remainders form the self-adjoint
matrix $\bigl(\begin{smallmatrix}R_o&R^{\rm new}_{\tlaw,p}\\R^{{\rm new}*}_{\tlaw,p}&R_o
\end{smallmatrix}\bigr)\succeq-(e_o+e_{\rm sg})qI$. As $q=r^2$, the prime Gram with
the shifts $dr$ is positive semidefinite for
$d\ge s_\tlaw\Gamma_A+(e_o+e_{\rm sg})r=d_\tlaw$, all evaluated at $(1/p,p^{-1/2})$. The
function $d_\tlaw$ is nondecreasing in $(q,r)$ (\eqref{eq:signed-constants}), so
this is the first hypothesis of \cref{lem:record:asymmetric-gram}(b) with
$c_{\rm L}=c_{\rm R}=d_\tlaw$.

\emph{Supports of at least two primes.} On every exact support $S\subseteq\Blk$
with $|S|\ge2$ the restrictions of the old and of the signed block operators
have norms with sum at most $c_{\ge2}/b$, with $c_{\ge2}$ at $(1/b,b^{-1/2})$, and
$c_{\ge2}r<2<s_\tlaw\le d_\tlaw$ (\cref{prop:signed-bounds}(d)); the bound comes from the
prime centred norms $(\Gamma_A+e_or)r$ of the reference and $c_+r$ of the three
components, through \cref{lem:sector}(a) and \cref{rem:signed-mixture-sectors},
the record being a combination of three products with absolute coefficient sum
$1+\tlaw$. By monotonicity and $r_b\le r_X$, $c_{\ge2}/b\le c_{\ge2}(X)r_Xr_b$, which is
the second hypothesis. The bound $c_{\ge2}r<2$ of the proof of
\cref{prop:signed-bounds}(d) holds at every $(q,r)$ with $q\le2^{-24}$ and
$r\le2^{-12}$, hence at $(q_X,r_X)$, as $X\ge2^{24}$ gives
$r_X=1/\lfloor\sqrt X\rfloor\le2^{-12}$; so $c_{\ge2}(X)r_X<2<s_\tlaw\le d_\tlaw(X)$, an
inequality that \cref{cert:signed-local} also asserts at the $33$ piece starts.

\emph{The block.} \Cref{lem:record:asymmetric-gram}(b) with $H_\Blk=\bar H_X$
gives the block Gram with the sides
\[
 d_{\rm L}=d_{\rm R}=\bar H_X\max\bigl(d_\tlaw(X),c_{\ge2}(X)r_X\bigr)=\bar H_Xd_\tlaw(X);
\]
here the block has at most $n_X$ primes
(\cref{lem:record:local-domain}(ii)), and $n_Xe_0(X)<1$ is a gate asserted in
\cref{cert:signed-local} and again by the evaluation of the increments of (ii)
(\cref{cert:record:matching}). The record of $F_2$ is the reflected transpose of
that of $F_1$, taken componentwise (\cref{def:record:mixture}(d)), so
\cref{lem:record:asymmetric-gram}(c) gives its block Gram with the same, equal,
sides. The Gram is an exact relation between the operators of the signed record
and of the reference, together with the operator bounds above; the total
variation $1+\tlaw$ of the record enters only the variation factor of (CP4)
(\cref{def:record:comparison-law}, \cref{rem:signed-mixture-sectors}), and the
Gram acquires no total-variation multiplier.

(ii) Since $\operatorname{up}(x)\ge x$ and the certified upper bound is at least
$\bar H_Xd_\tlaw(X)$, $\delta^+_{s,X}\ge\max(0,\bar H_Xd_\tlaw(X)-d^{\rm aux}_s(X))$, so
$d_s(X)\ge\max(\bar H_Xd_\tlaw(X),d^{\rm aux}_s(X))$, and by (i) and
\cref{lem:record:asymmetric-gram}(a) the block Gram holds with the sides $d_s(X)$,
for $F_1$ and, interchanged, for $F_2$. The identity for $k_{s,i}$ is the
definition, and $k_{s,i}\ge k^{\rm aux}_{s,i}\ge0$ because
$w_{{\rm D},i},\delta^+_{s,i}\ge0$.

(iii) By (ii), $(k_{{\rm L},i}+k_{{\rm R},i})\bar I_{2,i}=(k^{\rm aux}_{{\rm L},i}
+k^{\rm aux}_{{\rm R},i})\bar I_{2,i}+2w_{{\rm D},i}\cdot\frac12(\delta^+_{{\rm L},i}
+\delta^+_{{\rm R},i})\cdot\bar I_{2,i}$; bound each term by its rounding upwards
($w_{{\rm D},i},\bar I_{2,i}\ge0$) and sum over the pieces.
\end{proof}

At the first piece start $X=Q_3$: $n_X=34$, $\bar H_Xd_\tlaw(X)$ is about
$2.39034$, the auxiliary sides are about $2.05245$ and $2.18441$, and the
increments about $0.33789$ and $0.20593$. In every piece $\bar H_Xd_\tlaw(X)$
exceeds both auxiliary sides, so both increments are positive and both chosen sides are the auxiliary sides plus the rounded increments of \eqref{eq:signed-increment}, hence at least $\bar H_Xd_\tlaw(X)$; the auxiliary sides only
split the price $T'_K$ into the two certified sums of (iii): $T^{\rm aux\prime}_K$ is
about $0.0354169$, and the repriced increment price $tT^{\rm inc\prime}_K$ about
$0.1273584$. The coefficients
$k^{\rm aux}_{s,i}$ are positive in every piece, the smallest, at the last piece,
being about $0.0707$ (\cref{tab:data:sides}, \cref{cert:record:matching}).

\subsection{The dominators $M_{\rm L},M_{\rm R}$ and the total $T$}\label{subsec:matching-envelopes}

In every middle piece $i$ we use the chosen sides and the side coefficients
$k_{s,i}$ of \cref{prop:signed-sides}. For a middle block $\Blk$ of the piece:
\begin{itemize}
\item the \emph{chosen side diagonals} are
 \[
  K_{{\rm L},\Blk}=r_\Blk^2\bigl(k_{{\rm L},i}D^1_\Blk+k_{{\rm R},i}D^2_\Blk\bigr),\qquad
  K_{{\rm R},\Blk}=r_\Blk^2\bigl(k_{{\rm R},i}D^1_\Blk+k_{{\rm L},i}D^2_\Blk\bigr),
 \]
 without the normalization $1/\widetilde\mu_\Blk(1)$ that the companions in $\Nbr$
 carry; omitting it adds a nonnegative multiple of the same positive reserved
 diagonals;
\item the \emph{chosen matching diagonal} is
 $Q_{E,\Blk}=r_\Blk^2q_i\bigl(D_{E_0,\Blk}+D_{E_1,\Blk}\bigr)$, where
 $q_i=w_{{\rm S},i}(s_{E,i}+a_{E,i})$, and $s_{E,i}$, $a_{E,i}$ are the symmetric
 and skew coefficients of the matching records recorded in the native source
 (\cref{cert:record:means}): certified upper bounds for the coefficients
 $s_E(X_i)$, $a_E(X_i)$ of \cref{lem:record:exact-marginal-e}(ii) evaluated at $r_X=X_i^{-1/2}$, the least value allowed by (ii), which permits every $r_X\ge X_i^{-1/2}$ since the coefficients are nondecreasing in $(q,r)$ by (i); and
 $D_{E_s,\Blk}=\operatorname{Diag}_\iota(m_{E_s,\Blk})$ is the selected diagonal of $E_s$
 (\cref{def:record:eligible}), the diagonal copy at $\iota$ of the product law
 placed by $E_s$. By \cref{lem:record:exact-marginal-e}(ii) the symmetric and skew
 parts of the change of each matching record of $\Blk$, of weight $w_{\rm S}r_\Blk$
 normalized by $\widetilde\mu_\Blk(1)\ge1$, have norm envelopes
 $\tilde\pi^\epsilon D_{E_s,\Blk}$ with $\tilde\pi^{\rm s}+\tilde\pi^{\rm a}\le r_\Blk^2q_i$
 (\cref{subsec:variance:eligible}); so $Q_{E,\Blk}$ dominates the norm envelopes of
 the two matching changes, and replacing them by $Q_{E,\Blk}$ adds a positive
 reserved-diagonal measure. (The prices $S_{E,i}$, $A_{E,i}$ are formed with the
 coefficients of (ii) at $r_X=1/\lfloor\sqrt{X_i}\rfloor$, which are not smaller.)
\end{itemize}
Outside the middle, $K_{s,\Blk}=Q_{E,\Blk}=0$.

\paragraph{The matching exponent.}
With the totals at the exponent $a$ of \cref{def:closing-quantities} put
\[
 T_{QE}=S_E+A_E,\qquad
 T=S_E+A_E+S_{\rm oth}+A_{\rm oth}+T_K+S_{\rm u}(a)+A_{\rm u}(a)+10^{-12},
\]
with the $T_K$ of \cref{prop:signed-sides}(iii), the totals $S_E$, $A_E$ of the
matching records and $S_{\rm oth}$, $A_{\rm oth}$ of the other middle records
(\cref{lem:middle-comparison-prices}), and the upper and tail prices
$S_{\rm u}(a)$, $A_{\rm u}(a)$ of \cref{prop:record:upper}(ii).
\begin{itemize}
\item The gate $2q_i\bar I_{2,i}\le A_{E,i}+S_{E,i}$, with the matching prices of
 the piece at $a_0$ (\cref{def:closing-quantities}(3)), holds in every piece
 (\cref{cert:record:means}), so the chosen matching diagonals cost at most
 $T_{QE}$.
\item The prices of the prefix records lie inside the allowance $10^{-12}$
 (\cref{lem:frontier:iidenergy}(ii)--(iv)).
\item The symmetric and skew prices of the fold records do not enter $T$; only
 $T_K$ does.
\end{itemize}
At the baseline exponent $T_K$ is about $0.1628311$ and $T_{QE}$ about $0.4233528$
(\cref{cert:closing}).

\begin{proposition}[Matching envelopes and the matching exponent]
\label{prop:record:matching}
\leavevmode
\begin{enumerate}[label=(\roman*)]
\item For every atom $\Blk$ there are measures $M_{{\rm L},\Blk}$ and
 $M_{{\rm R},\Blk}$ with $\bigl(\begin{smallmatrix}M_{{\rm L},\Blk}&\nu_\Blk\\
 \nu_\Blk^*&M_{{\rm R},\Blk}\end{smallmatrix}\bigr)\succeq_\iota0$, such that
 $M_{{\rm L},\Blk}-\Nbr_\Blk$ and $M_{{\rm R},\Blk}-\Nbr_\Blk$ are positive
 reserved-diagonal measures whose normalizations have uniform one-slot
 marginals and do not depend on the indices. In a middle block
 \[
  M_{s,\Blk}=\Nbr_\Blk+Q^\circ_\Blk+Q_{E,\Blk}+K_{s,\Blk},\qquad s\in\{{\rm L},{\rm R}\},
 \]
 where $Q^\circ_\Blk$ is the sum of the norm envelopes of
 \cref{def:record:comparison-law} of every record change other than the two
 fold records and the two matching records; outside the middle
 $M_{s,\Blk}=\mu_\Blk+Q_\Blk$, with $Q_\Blk$ the sum of the norm envelopes of the
 changes at $\Blk$.
\item For both $s$, with the $T$ of the paragraph above,
 \[
  \sum_\Blk\operatorname{price}(M_{s,\Blk}-\Nbr_\Blk)\le T,\qquad
  T\in[0.607496889410,\,0.607496889411]\quad\text{at }a=\PushA.
 \]
 So (CP6) holds with the matching exponent $T$ of
 \cref{def:closing-quantities}(5).
\item The single-half random-inclusion energies $V_{\rm L}$ of $X$ under
 $\bigotimes_\Blk M_{{\rm L},\Blk}$ and $V_{\rm R}$ of $Y$ under
 $\bigotimes_\Blk M_{{\rm R},\Blk}$ (\cref{def:hyb:coupled}) are at most
 $(2^6J+o(1))Y$, with $J=Be^T(R_I/2)^6$.
\end{enumerate}
\end{proposition}

\begin{proof}
(i) In a middle block, by \cref{def:record:frame}(d),
$\Nbr_{{\rm res},\Blk}-\lambda_{{\rm res},\Blk}=H_\Blk\succeq_\iota0$ and
\[
 \Nbr_{{\rm res},\Blk}+\lambda_{{\rm res},\Blk}=\bigl(2\mu_\Blk-2F^{(2)}_\Blk
 -2F_{{\rm fold},\Blk}-2E_{{\rm fold},\Blk}+H_\Blk\bigr)+2E_{{\rm new},\Blk}+2G_\Blk .
\]
The proof of \cref{lem:record:gram}(ii) writes
$2\mu_\Blk-F_{{\rm fold},\Blk}-2E_{{\rm fold},\Blk}+H_\Blk$ as a sum of measures each
positive semidefinite at $\iota$, one for every component of $\widetilde\mu_\Blk$
(all normalized by $\widetilde\mu_\Blk(1)$): for each fold record $F_j$, $j=1,2$,
the pair $2w_{\rm D}r\Pi_{F_j,\Blk}+v_{\rm D}r^2D^j_\Blk$, positive at $\iota$ by the
protection \eqref{eq:protection} (\cref{lem:front12:reference}(iv)); for each
matching record the record with its part of $H_\Blk$; and every other component
together with its own one-layer companion. The two fold pairs add up to
exactly $2F^{(2)}_\Blk+F_{{\rm fold},\Blk}$. Removing them leaves a positive
semidefinite sum, and $E_{\rm new}$ and $G$ are positive and reserved diagonal.
Since $\lambda_{{\rm res},\Blk}$ is symmetric at $\iota$, the symmetric-bracket
criterion of \cref{lem:record:block-gram}(d) gives
$\bigl(\begin{smallmatrix}\Nbr_{\rm res}&\lambda_{\rm res}\\
\lambda_{\rm res}&\Nbr_{\rm res}\end{smallmatrix}\bigr)\succeq_\iota0$; no
diagonal is used twice.

Add $w_{\rm D}r/\widetilde\mu_\Blk(1)$ times the block Grams of
\cref{prop:signed-sides}(i),(ii) for the signed $F_1$ record, with the chosen
sides $d_{\rm L},d_{\rm R}$, and, with the sides interchanged, for the signed $F_2$
record. The off-diagonal entry becomes
$\lambda_{\rm res}+w_{\rm D}r(\Pi'_{F_1}+\Pi'_{F_2})/\widetilde\mu_\Blk(1)$, which is
$\lambda_\Blk$ with the two fold records replaced by their signed comparison
records. On the layer-zero side the diagonal entry becomes
\[
 \Nbr_{\rm res}+F^{(2)}_\Blk+\frac{w_{\rm D}r^2}{\widetilde\mu_\Blk(1)}\bigl(d_{\rm
 L}D^1_\Blk+d_{\rm R}D^2_\Blk\bigr)
 =\Nbr_\Blk+\frac{r^2}{\widetilde\mu_\Blk(1)}\bigl(k_{\rm L}D^1_\Blk+k_{\rm R}D^2_\Blk\bigr),
\]
since $\Nbr_{\rm res}+F^{(2)}_\Blk=\Nbr_\Blk-F_{{\rm fold},\Blk}$ and
$F_{{\rm fold},\Blk}=v_{\rm D}r^2(D^1_\Blk+D^2_\Blk)/\widetilde\mu_\Blk(1)$; the difference between $\Nbr_\Blk+K_{{\rm L},\Blk}$ and this diagonal entry is a positive reserved-diagonal measure,
because $k_s\ge0$ and $\widetilde\mu_\Blk(1)\ge1$. On the layer-one side it is at
most $\Nbr_\Blk+K_{{\rm R},\Blk}$ in the same way. Enlarging a diagonal entry by a
positive reserved-diagonal measure keeps the bracket. Finally add the brackets
of \cref{def:record:comparison-law} of every other record change (they are
those of \cref{lem:record:record-expansion}): for the two matching changes with
the diagonal $Q_{E,\Blk}$, which dominates their norm envelopes
(\cref{subsec:matching-envelopes}), and for the other changes with their norm
envelopes, whose sum is $Q^\circ_\Blk$. Their off-diagonal entries complete
$\nu_\Blk$. Outside the middle only this last step occurs, starting from
$\bigl(\begin{smallmatrix}\mu_\Blk&\mu_\Blk\\ \mu_\Blk&\mu_\Blk
\end{smallmatrix}\bigr)\succeq_\iota0$ ($\mu_\Blk$ is reflection positive at
$\iota$, \cref{hyp:record:reference}). The side gates make
$k_{\rm L},k_{\rm R}\ge0$; the added diagonals are positive, reserved diagonal,
with uniform one-slot marginals, and independent of the indices. The exchange
of \cref{lem:record:asymmetric-gram}(c) is used here, so that each envelope
carries $k_{\rm L}+k_{\rm R}$ for the two fold records together; without it one
side would carry $2k_{\rm R}$.

(ii) The price of $K_{s,\Blk}$ is its mass times $q_\Blk^{2a}$, at most
$(k_{{\rm L},i}+k_{{\rm R},i})r_\Blk^2q_\Blk^{2a}$, and $q_\Blk^{2a}\le tq_\Blk^{2a_0}$ for a
middle block (\cref{lem:record:repricing}; $\log q_\Blk\le L$). Summed over the
blocks of the piece $i$ the chosen side diagonals cost at most
$t(k_{{\rm L},i}+k_{{\rm R},i})\bar I_{2,i}$, and over all pieces at most $tT'_K=T_K$
by \cref{prop:signed-sides}(iii). In the same way $Q_{E,\Blk}$ has the price
$2q_ir_\Blk^2q_\Blk^{2a}$, so the chosen matching diagonals cost at most
$t\sum_i2q_i\bar I_{2,i}\le t\sum_i(A_{E,i}+S_{E,i})=T_{QE}$ by the gate; this is
the budget of the norm envelopes they replace, and no second copy is
introduced. The norm envelopes $Q^\circ_\Blk$ of the merges and of the deleted
records have the prices of \cref{lem:middle-comparison-prices}, of total at most
$S_{\rm oth}+A_{\rm oth}$ at $a$. Outside the middle, $Q_\Blk$ is the norm envelope
of the upper and tail repairs and deletions, of total price at most
$S_{\rm u}(a)+A_{\rm u}(a)$ (\cref{prop:record:upper}(ii)), and of the deleted
exception of the prefix, charged to the symmetric part only, of total price
below $10^{-12}$ (\cref{lem:frontier:iidenergy}(ii)--(iv) with $\mcl=8192$). Adding
these bounds gives $T$; its enclosure is that of \cref{cert:closing}.

(iii) Each $M_{s,\Blk}-\Nbr_\Blk$ is a nonnegative combination of positive
reserved-diagonal laws with uniform one-slot marginals, independent of the
index, with $\sum_\Blk\operatorname{price}(M_{s,\Blk}-\Nbr_\Blk)\le T$ by (ii). \Cref{lem:envelope-energy} with
$Q_s=M_{s,\Blk}-\Nbr_\Blk$ and $T_Q=T$ gives $V_s\le(Be^TR_I^6+o(1))Y=
(2^6J+o(1))Y$, the conductor argument being carried out inside each
half-energy. No symmetry of the comparison records and no odd-skew
cancellation is used.
\end{proof}

\subsection{The common mean for the selected law}\label{subsec:cm-selected}

For the final construction the six queries $\mathcal E_6$ have pairwise
disjoint endpoints (\cref{lem:slot-array}(iv)) and the index law, with
$\delta_j=1$ for the interval coordinate, for the $2$-adic coordinate modulo $8$, and for the prime $3$ coordinate, and
$\delta_p=\sqrt{p+1}/(p-1)$ at the class-three primes $7\le p<Q_3$, satisfies
(CP0)(d) (\cref{lem:imp:padic-envelope}); the matching envelopes of
\cref{prop:record:matching}(i) are brackets as in (CP6). So
\cref{cor:index-support-frame} applies to the signed index form $J_\nu$ of
\eqref{eq:record:expansion}. With the \emph{mean energies}
$e_s=Q_{\otimes M_s}(\bar X,\bar X)$ of the two envelopes ($\bar X$ the
half-function with $F_A$ at every slot) and the single-half random-inclusion
energies $V_{\rm L},V_{\rm R}$ of \cref{prop:record:matching}(iii), it gives
\begin{equation}\label{eq:record:cm-selected}
 |J_\nu-Y_\nu|\le\tfrac12\bigl[(V_{\rm L}-e_{\rm L})+(V_{\rm R}-e_{\rm
 R})\bigr],\qquad V_{\rm L},V_{\rm R}\le(2^6J+o(1))Y .
\end{equation}
The empty-support components are kept exactly: on one side $Y_\nu$, on the
other the mean energies, which exceed the mean energy $Y_\nu$ of the law
itself. This saves one of the $64$ copies of the reserved cost in $\Theta$
(\cref{thm:record:criterion}) and is why the lower bounds of the next
subsection enter $\Theta$.

\subsection{Mean energies exceed $mY$ with $m>1$}\label{subsec:record:positive_means}

The number $\Theta$ in \cref{thm:comparison} is nonincreasing
in the mean credit $m$ of (CP7): a larger proved $m$ only helps.
We bound
the mean energies from the reference side. Each envelope contains $\Nbr$,
and $\Nbrz$ is $\mu$ plus positive reserved-diagonal records; at $\iota$
every monomial of the product of $\Nbrz$ and the chosen diagonals is
therefore positive semidefinite, and explicit lower bounds for a few
monomials may be kept while all others are discarded. The monomials kept are
those of the polarized fold diagonals, the chosen side diagonals and the
chosen matching diagonals, on sets of at most four distinct blocks, and of
five or six blocks of a pool of late pieces. Their energies are bounded below
by restoration (\cref{lem:record:restoration}) at the chamber cut where the
two-layer fold record is the diagonal copy of its own induced law, and, for the
matching diagonals, by the aligned bridge (\cref{lem:record:paired-e-joint}). At
the baseline exponent these coupons are necessary for
$\Theta<1$ in this construction: they lower $\Theta$ by about
$(m-1)/64\approx0.0092$, against a margin $1-\Theta$ of about
$2.67\cdot10^{-11}$.

\paragraph{Notation.}
Energies are those of the toolkit
\ref{it:record:E1}--\ref{it:record:E3}: $Y[\kappa]=Y_{\kappa\otimes\mu_{-X}}$
for a measure $\kappa$ on the labels of a set $X$ of atoms, and
$Y^{\Nbr}[\kappa]$, $Y^{\Nbrz}[\kappa]$ with $\Nbr_{-X}$, $\mathsf N_{0,-X}$
outside. Let $\Blk$ be a middle block of the piece $[X_i,Y_i)$
(\cref{def:record:exact-blocks}), with least prime $b_\Blk$ and
$r_\Blk=b_\Blk^{-1/2}$. For $j\in\{1,2\}$ the normalized one-layer
companion $D^j_\Blk$ of $F_j$ is a positive normalized law, fixed samplewise
by $\iota$, with uniform one-slot marginals, and so are the selected diagonals
$D_{E_s,\Blk}$ of the matching records; $D_{E_s,\Blk}=\operatorname{Diag}_\iota(\Pi^{(1)}_{E_s})$
is the one-layer matching law on the reserved diagonal, whereas $D^j_\Blk$ is the diagonal copy of $\nu_F^{\otimes\Blk}$ placed by $F_j$, and the selected diagonal $D_{F_j,\Blk}=\operatorname{Diag}_\iota(m^{F_j}_{0,\Blk})$ is obtained from it by the density of $m^{F_j}_{0,\Blk}$ relative to $\nu_F^{\otimes\Blk}$. By \cref{def:record:frame} and \cref{lem:record:gram},
\[
 (\Nbrz-\mu)_\Blk=\mathsf E_\Blk+H_\Blk,\qquad\mathsf E_\Blk=(E_{\rm add}+G)_\Blk,
 \qquad E_{{\rm add},\Blk}=\frac{v_{\rm D}r_\Blk^2}{1+h_\Blk}\bigl(D_{E_0,\Blk}+D_{E_1,\Blk}\bigr),
\]
with $\widetilde\mu_\Blk(1)=1+h_\Blk$, where $h_\Blk$ is the nonroot mass
\eqref{eq:record:nonroot-mass}; all summands are nonnegative and reserved
diagonal. The fold companion
$F_{{\rm fold},\Blk}=v_{\rm D}r_\Blk^2(1+h_\Blk)^{-1}(D^1_\Blk+D^2_\Blk)$ is a summand of
$\Nbr$, not of $\Nbrz$; it enters the mean only through the polarization orders
$\mathsf E_\Blk\pm F_{{\rm fold},\Blk}\succeq_{\vartheta_0}0$
(\cref{lem:middle-polarization}, \cref{rem:polarization-uses}(3)). By
\cref{prop:record:matching}(i),
$M_{s,\Blk}-\Nbr_\Blk-K_{s,\Blk}-Q_{E,\Blk}=Q^\circ_\Blk\succeq_\iota0$
(\ref{it:record:E3}).

For the piece $i$ we use the piece values $h_i$, $M_i$, $\chi_i$,
$\varepsilon_i$ and the density factors $H_i=(1+2/(X_i-4))^{n_i}$, with
$n_i=\lfloor L/\log_-X_i\rfloor$ and $L=570$, of \eqref{eq:variance:pieces}, and
the caps $C_i$ of \cref{lem:piece-caps}. For a multiset $m$ of pieces, $C_m$,
$\eta_m$ and the restored value $\ell_m=\mathcal F(C_m,h_m,M_m,\chi_m)$ are those
of the piece data of \cref{subsec:variance:backgrounds}, and a set of distinct
middle blocks has the \emph{pattern} $m$ if $m_i$ of its blocks lie in the
piece $i$. (Multisets of pieces are written $m$, as in \cref{subsec:variance:backgrounds}, and ``pattern'' here means such a multiset, not a chamber pattern $\phi$; they occur only as indices and in the
sums over $|m|$, and are not to be confused with the mean credit $m$ of
(CP7).) The factor $H_i$ bounds the density of the actual block marginal $m_{0,\Blk}$
relative to $\nu_F^{\otimes\Blk}$ for every block of the piece: at a prime,
$\sup g_0=(p-2)/(p-4)=1+2/(p-4)$, the value of the formula of
\cref{lem:old-f-law}(a) at $e_{01}=e_{12}=0$, and a block of the piece has at most
$n_i$ primes, all at least $X_i$ (\cref{lem:record:local-domain}). Put
\[
 c_{E,i}=\frac{v_{{\rm D},i}}{1+h_i},\qquad c_{F,i}=\frac{c_{E,i}}{H_i},
\]
the coefficients of the matching diagonals of $E_{\rm add}$ and of the polarized
fold diagonals, and let $W_0(m)=W(\varnothing,m;\mathbf1)$ be the distinct-choice
weight of \cref{lem:distinct-choice}(a) with unit coefficients: the sum of
$\prod_{\Blk\in X}b_\Blk^{-1}$ over the sets $X$ of distinct middle blocks of
pattern $m$ is at least $W_0(m)$. Except through the caps $C_i$, these
quantities do not depend on $a$.

\begin{lemma}[Restored companion energies]\label{lem:record:companion-energy}
Let $X$ be a nonempty set of distinct middle blocks, $\Blk\in X$ lying in the
piece $i(\Blk)$, and $j\in\{1,2\}$.
\begin{enumerate}[label=(\alph*)]
\item For every real $\ell$, if $Y[\bigotimes_{\Blk\in X}\Pi_{F_j,\Blk}]\ge(\ell-o(1))Y$, then
 $Y[\bigotimes_{\Blk\in X}D^j_\Blk]\ge\bigl(\ell/\prod_{\Blk\in X}H_{i(\Blk)}-o(1)\bigr)Y$.
 In particular, if $X$ has the pattern $m$, then
 $Y[\bigotimes_{\Blk\in X}D^j_\Blk]\ge(\ell_m/\prod_{i\in m}H_i-o(1))Y$.
\item If $X$ has the pattern $m$, then
 $Y[\bigotimes_{\Blk\in X}D_{E_0,\Blk}]+Y[\bigotimes_{\Blk\in X}D_{E_1,\Blk}]
 \ge2\bigl([\ell_m-C_m\eta_m]_+-o(1)\bigr)Y$.
\end{enumerate}
\end{lemma}

\begin{proof}
(a) Put $\Pi_X=\bigotimes_{\Blk\in X}\Pi_{F_j,\Blk}$.

\emph{Contraction at the reserved cut.} By
\cref{lem:record:fold-identification}(d), $\Pi_X$ is symmetric at $\iota$
with actual half marginal $m_X=\bigotimes_{\Blk\in X}m^{F_j}_{0,\Blk}$, so
the contraction \ref{it:record:P1} gives
$Y[\Pi_X]\le Y[\operatorname{Diag}_\iota(m_X)]$.

\emph{Density.} For a block $\Blk$ of the piece $i$, the density of
$m^{F_j}_{0,\Blk}$ with respect to the half law $\nu_F^{\otimes\Blk}$ (placed by
$F_j$) of $D^j_\Blk$ is at most $H_i$ (paragraph above), and densities of
products multiply. So the density domination \ref{it:record:P2}, a domination
of reserved-diagonal measures, gives
$Y[\operatorname{Diag}_\iota(m_X)]\le\prod_{\Blk\in X}H_{i(\Blk)}\,
Y[\bigotimes_{\Blk\in X}D^j_\Blk]$. Combining the two steps with the hypothesis
proves the first claim. For the second, the blocks of $X$ share the fold chamber
$\vartheta_{2j}$ of $F_j$ (\cref{lem:record:fold-identification}(b)), and the piece
data give $Y[\Pi_X]\ge(\ell_m-o(1))Y$ (\cref{subsec:variance:backgrounds},
\cref{lem:record:restoration}, with the cap $C_m\ge C_X$ of \cref{lem:piece-caps}).

(b) By \cref{subsec:variance:backgrounds}, \cref{lem:record:paired-e-joint}(b)
holds for $X$ with $l_0=\ell_m$, $C=C_m$ and $\varepsilon_X\le\eta_m$; and the
actual half marginal at $\iota$ of the aligned embedding $E^X_s$ is
$\bigotimes_{\Blk\in X}m_{E_s,\Blk}$, so its diagonal copy is
$\bigotimes_{\Blk\in X}D_{E_s,\Blk}$. Hence (b) is \cref{lem:record:paired-e-joint}(c),
as $[l_0-C\varepsilon_X]_+\ge[\ell_m-C_m\eta_m]_+$.
\end{proof}

\begin{proposition}[Mean credit]\label{prop:record:mean-credit}
Assume, in every middle piece, the protection \eqref{eq:protection}, the side
gates $k^{\rm aux}_{s,i}\ge0$ (\cref{prop:signed-sides}) and the matching gates
$2q_i\bar I_{2,i}\le A_{E,i}+S_{E,i}$ (\cref{subsec:matching-envelopes}). Then, for
$s\in\{{\rm L},{\rm R}\}$, $e_s\ge(m-o(1))Y$ with
\begin{equation}\label{eq:mean-credit}
 m=1+m_F+m_Q-R_{\rm err}e^{T_K+T_{QE}},
\end{equation}
where the sums run over the multisets $m$ of pieces, the products over the
pieces of $m$ with multiplicity,
\begin{align*}
 m_F&=\sum_{1\le|m|\le4}W_0(m)\,\ell_m\Bigl[2\prod_{i\in m}c_{F,i}
 +\prod_{i\in m}\frac{k_{{\rm L},i}}{H_i}+\prod_{i\in m}\frac{k_{{\rm R},i}}{H_i}\Bigr]
 +m^{\rm pool}_F,\\
 m_Q&=\sum_{1\le|m|\le4}2W_0(m)\,[\ell_m-C_m\eta_m]_+
 \Bigl[\prod_{i\in m}(c_{E,i}+q_i)-\prod_{i\in m}c_{E,i}\Bigr]+m^{\rm pool}_Q,
\end{align*}
and $m^{\rm pool}_F$, $m^{\rm pool}_Q$ are the late-pool coupons of degrees five and
six of \cref{lem:pool-coupons}(b), on the pieces starting at $2^{31}$ or later.
Hence (CP7) holds with the mean credit $m$.
\end{proposition}

\begin{proof}
\emph{Step 1: from the envelopes to $\Nbrz$.} In every middle block
$0\preceq_\iota\Nbr_\Blk+K_{s,\Blk}+Q_{E,\Blk}\preceq_\iota M_{s,\Blk}$: the bracket
measure is positive semidefinite at $\iota$ (\cref{lem:record:gram}(ii)), the
chosen diagonals are positive and reserved diagonal (their coefficients are
nonnegative by the gates), and the difference is $Q^\circ_\Blk\succeq_\iota0$;
outside the middle $M_{s,\Blk}-\mu_\Blk=Q_\Blk\succeq_\iota0$. By the tensor
monotonicity of \cref{lem:record:block-gram}(d), applied to the cut matrices at
$\iota$ \eqref{eq:cc:kron}, $\bigotimes_\Blk M_{s,\Blk}\succeq_\iota
\bigotimes_\Blk(\Nbr_\Blk+K_{s,\Blk}+Q_{E,\Blk})$, and evaluation at $\bar X$ gives
$e_s\ge Y_{\widetilde M_s}$ with $\widetilde M_s=\bigotimes_\Blk(\Nbr_\Blk+K_{s,\Blk}+Q_{E,\Blk})$.
Let $M_{0,s}$ be the same product with $\Nbrz$ in place of $\Nbr$. Expand both
over the set $\mathcal S$ of blocks that carry $K_{s,\Blk}$ or $Q_{E,\Blk}$ and
over, at each block of $\mathcal S$, the choice of $K_{s,\Blk}$ or $Q_{E,\Blk}$. Each resulting product $C_{\mathcal S}$ of chosen diagonals (the notation suppresses the choice) is a positive measure whose normalization
has uniform one-slot marginals and is admissible (\cref{prop:record:matching}(i)),
and its price is the product of the prices $\pi_\Blk$ of its factors; by
\ref{it:record:E1},
$|Y^{\Nbr}[C_{\mathcal S}]-Y^{\Nbrz}[C_{\mathcal S}]|\le(1+o(1))R_{\rm err}
\prod_{\Blk\in\mathcal S}\pi_\Blk Y$. Summing over all choices gives at most
$(1+o(1))R_{\rm err}\prod_\Blk\bigl(1+\pi(K_{s,\Blk})+\pi(Q_{E,\Blk})\bigr)Y\le
(1+o(1))R_{\rm err}e^{T_K+T_{QE}}Y$, because the chosen side diagonals cost at most
$T_K$ and the chosen matching diagonals at most $T_{QE}$, for either $s$ (proof of
\cref{prop:record:matching}(ii)). So
$e_s\ge Y_{M_{0,s}}-(R_{\rm err}e^{T_K+T_{QE}}+o(1))Y$.

\emph{Step 2: the expansion.} Write
$M_{0,s}=\bigotimes_\Blk\bigl(\mu_\Blk+\mathsf E_\Blk+H_\Blk+K_{s,\Blk}+Q_{E,\Blk}\bigr)$ and
expand $\mathsf E_\Blk=E_{{\rm add},\Blk}+G_\Blk$ as well. The cut matrix at $\iota$ of
every term is a tensor product of positive semidefinite matrices, since
$\mu_\Blk\succeq_\iota0$ and the other measures are positive and reserved diagonal;
so every term has nonnegative energy (\ref{it:record:E2}), and we may keep
some and discard the rest. Keep the term with $\mu$ at every atom, whose
energy is $Y$, and, for every set $X$ of distinct middle blocks with
$1\le|X|\le4$, or with $|X|\in\{5,6\}$ and all blocks in pieces of the pool, the
three families of terms in which the blocks outside $X$ carry $\mu$ and the
blocks of $X$ carry
\begin{itemize}
\item[($\alpha$)] labels in $\{E_{\rm add},G\}$, together
 $\bigotimes_{\Blk\in X}\mathsf E_\Blk$;
\item[($\beta$)] the chosen side diagonal, together $\bigotimes_{\Blk\in X}K_{s,\Blk}$;
\item[($\gamma$)] labels in $\{E_{\rm add},Q_E\}$, at least one of them $Q_E$.
\end{itemize}
These are distinct terms of the expansion: ($\gamma$) has a $Q_E$ label, and
($\alpha$), ($\beta$) have none and differ in every label. All other terms are
discarded. The number of middle blocks is finite and does not depend on the
scale, so the errors $o(1)Y$ below add up to $o(1)Y$.

\emph{Step 3: the empty-$Q$ branch.} ($\alpha$) By the polarization orders
$\mathsf E_\Blk\pm F_{{\rm fold},\Blk}\succeq_{\vartheta_0}0$
(\cref{rem:polarization-uses}(3)) and the identity of all degrees
\cref{lem:tr:polarization}(d),
\[
 \bigotimes_{\Blk\in X}\mathsf E_\Blk-\bigotimes_{\Blk\in X}F_{{\rm fold},\Blk}
 =2^{-|X|}\sum_{\sigma\in\{\pm1\}^X}\Bigl(1-\prod_\Blk\sigma_\Blk\Bigr)
 \bigotimes_{\Blk\in X}\bigl(\mathsf E_\Blk+\sigma_\Blk F_{{\rm fold},\Blk}\bigr)
 \succeq_{\vartheta_0}0 .
\]
The factors $\mu_{\Blk'}$, $\Blk'\notin X$, are reflection positive at
$\vartheta_0$, so $Y[\bigotimes\mathsf E_\Blk]\ge Y[\bigotimes F_{{\rm fold},\Blk}]$
(\ref{it:record:E2}). Expand $\bigotimes F_{{\rm fold},\Blk}$ in the measures
$\bigotimes_\Blk D^{j_\Blk}_\Blk$; the products with mixed orientations are positive
and reserved diagonal, with nonnegative energy, and are discarded; the two
common orientations $j\equiv1$, $j\equiv2$ have, by
\cref{lem:record:companion-energy}(a), energies at least $\ell_m/\prod_{i\in m}H_i-o(1)$
times $Y$, with the coefficient $\prod_\Blk v_{\rm D}r_\Blk^2/(1+h_\Blk)\ge
\prod_\Blk r_\Blk^2\prod_{i\in m}c_{E,i}$. So the terms ($\alpha$) of $X$ have energy at
least $(2\prod_\Blk r_\Blk^2\,\ell_m\prod_{i\in m}c_{F,i}-o(1))Y$, the factor $2$ counting
the two common orientations. The fold companions are summands of $\Nbr$ and not
of $\Nbrz$, so this polarized copy is the only way in which they enter the
mean.

($\beta$) Expand $\bigotimes K_{s,\Blk}$ in the same measures and keep the two common
orientations: for $s={\rm L}$ they are
$\prod_\Blk r_\Blk^2\bigl(\prod_\Blk k_{{\rm L},i(\Blk)}\bigotimes D^1_\Blk
+\prod_\Blk k_{{\rm R},i(\Blk)}\bigotimes D^2_\Blk\bigr)$, and for $s={\rm R}$ the same with
$D^1$ and $D^2$ exchanged. By \cref{lem:record:companion-energy}(a) both give at
least $(\prod_\Blk r_\Blk^2\,\ell_m[\prod_{i\in m}k_{{\rm L},i}/H_i+\prod_{i\in m}k_{{\rm R},i}/H_i]
-o(1))Y$: the exact left and right pairing is retained, no ordering of the
sides is needed, and both envelopes receive the same credit. The chosen side
diagonals carry no factor $1/(1+h_\Blk)$, so the side ratios are $k_{s,i}/H_i$. In
the bracket of one fold copy the measure $\Nbrz-\mu$ contains no fold diagonal
(\cref{lem:record:gram}), so there are no terms mixing fold diagonals of $\Nbrz$
with the chosen side diagonals.

\emph{Step 4: the nonempty-$Q$ branch.} ($\gamma$) Every factor is
$c_\Blk r_\Blk^2(D_{E_0,\Blk}+D_{E_1,\Blk})$ with
$c_\Blk\in\{v_{\rm D}/(1+h_\Blk),\,q_{i(\Blk)}\}$, so the terms ($\gamma$) of $X$ add up
to
\[
 \prod_\Blk r_\Blk^2\Bigl[\prod_\Blk\Bigl(\frac{v_{\rm D}}{1+h_\Blk}+q_{i(\Blk)}\Bigr)
 -\prod_\Blk\frac{v_{\rm D}}{1+h_\Blk}\Bigr]\bigotimes_\Blk\bigl(D_{E_0,\Blk}+D_{E_1,\Blk}\bigr).
\]
The bracket is a polynomial with
nonnegative coefficients that is nondecreasing in each $v_{\rm D}/(1+h_\Blk)\ge
c_{E,i(\Blk)}$, so it is at least $\prod_{i\in m}(c_{E,i}+q_i)-\prod_{i\in m}c_{E,i}$.
Keep the two aligned orientations of $\bigotimes(D_{E_0,\Blk}+D_{E_1,\Blk})$ and discard
the others; by \cref{lem:record:companion-energy}(b) their energies add up to at
least $2([\ell_m-C_m\eta_m]_+-o(1))Y$. The empty and nonempty $Q$ branches are
disjoint, even where their physical diagonals are proportional.

\emph{Step 5: summation.} For $1\le|X|\le4$ the bounds of Steps 3 and 4 depend on
$X$ only through its pattern $m$ and the factor $\prod_\Blk r_\Blk^2$, and the sum of
this factor over the sets of pattern $m$ is at least $W_0(m)$
(\cref{lem:distinct-choice}(a)); all coefficients are nonnegative. This gives the
sums over $1\le|m|\le4$. For a set $X$ of $k\in\{5,6\}$ blocks of pieces of the
pool, the first assertion of \cref{lem:pool-coupons} gives
$Y[\bigotimes_{\Blk\in X}\Pi_{F_j,\Blk}]\ge(\ell_{\mathcal P^k}-o(1))Y$ and
\cref{lem:record:paired-e-joint}(b), (c) with the pool data, so Steps 3 and 4 hold
with $\ell_{\mathcal P^k}$, $C_{\mathcal P}^k$ and $\eta_{\mathcal P^k}$ in place of
$\ell_m$, $C_m$ and $\eta_m$ (\cref{lem:record:companion-energy}(a), and the proof of
\cref{lem:record:companion-energy}(b) with these data);
the bounds of \cref{lem:pool-coupons}(b) for the sums over these sets of the
coefficients of the three branches give the coupons $m^{\rm pool}_F$ and
$m^{\rm pool}_Q$, which occupy degrees disjoint from those of the other kept terms.
Adding the kept terms and Step 1 gives
$e_s\ge(1+m_F+m_Q-R_{\rm err}e^{T_K+T_{QE}}-o(1))Y$ for both $s$. Every restoration is recomputed at the exponent $a$ with the product caps $C_m$ at that exponent, not transported from an estimate at $a_0$, and no lower estimate
of an upper-priced norm is used.

The closing number $\Theta$ is nonincreasing in $m$ (\cref{rem:cp:monotone}), so
certified lower bounds for $m_F$, $m_Q$ and upper bounds for $R_{\rm err}$, $T_K$ and
$T_{QE}$ may be used.
\end{proof}

The chosen side diagonals carry no factor $1/(1+h_\Blk)$, so the side ratios
are $k_{s,i}/H_i$. The fold companions of $\Nbr$ carry it, so $c_{F,i}$ keeps
$1/(1+h_i)$. The full chosen prices $T_K$ and $T_{QE}$ are charged in the error
exponent, as in the matching exponent. The restorations are
evaluated at the product caps of the final exponent, with piece
data that do not depend on $a$. The following numerical values describe
the baseline certificate, before the joint allocation below; they are in
\cref{cert:record:means}:
$m_F>0.400172842378$, $m_Q>0.190134971388$, of which the terms
pooled over the pieces from $2^{31}$ contribute
about $1.1951\cdot10^{-7}$ and $6.5884\cdot10^{-7}$ (the degree-six parts of both pooled coupons $m^{\rm pool}_F$ and $m^{\rm pool}_Q$ vanish),
$R_{\rm err}e^{T_K+T_{QE}}<3.67071778633\cdot10^{-4}$ and $m>1.589940741988$.

\subsection{Allocating the full positive mean jointly with the comparison}
\label{subsec:joint-mean}

The closing inequality contains the comparison energy with coefficient
$63/64$ and the mean energy with coefficient $-1/64$. A positive
diagonal occurring in both can therefore reduce the closing bound:
subtract its mean coefficient from $63$ times its comparison coefficient,
then apply an upper bound to the remaining nonnegative energy.
This is the joint allocation of \cref{cor:joint-closing}.
We apply it to the positive expansion in Step~2 of
\cref{prop:record:mean-credit}. Of the lower bounds retained there
(called mean credits), only those supported on one block are used
separately here. The following argument replaces the credits on
two or more blocks.

For a middle block write
\[
 \alpha_B=\frac{w_{{\rm D},B}r_B^2}{1+h_B},\qquad
 T_B=E_{{\rm add},B}+H_B
     =\beta_B(D_{E_0,B}+D_{E_1,B}),\qquad
 \mathsf E_B=E_{{\rm add},B}+G_B.
\]
The measures $E_{{\rm add},B},G_B,H_B$, the fold diagonals $D_{F_j,B}$,
and the matching diagonals $D_{E_j,B}$ are those defined in
\cref{subsec:record:positive_means}; $K_{s,B}$ is the added diagonal
on side $s\in\{{\rm L},{\rm R}\}$ of the matching bracket
(\cref{prop:record:matching}). Here $\beta_B$ is a local coefficient.
The actual chosen diagonals satisfy
\begin{equation}\label{eq:joint-resources}
 \begin{aligned}
 2\le\beta_B/\alpha_B&\le R=2.004317,\qquad\sigma=0.3842,\\
 Q_{E,B}&\succeq_\iota\alpha_B(D_{E_0,B}+D_{E_1,B}),\\
 K_{s,B}&\succeq_\iota\sigma\alpha_B(D_{F_1,B}+D_{F_2,B}).
 \end{aligned}
\end{equation}
The density of an actual fold marginal relative to its normalized
companion is at most $H_*=1.000003310$. Thus its companion diagonal
dominates $H_*^{-1}D_{F_j,B}$.
These assertions use the same $33$ physical pieces and chosen sides
as \cref{prop:record:matching}. Explicitly, on a piece starting at $X$,
\[
 \frac{\beta_B}{\alpha_B}
 \le\frac{v_{\rm D}}{w_{\rm D}}
       +\frac{6}{\lfloor\sqrt X\rfloor(1-3/X)}\le R,
 \qquad
 \frac{\min(k_{{\rm L},i},k_{{\rm R},i})}{w_{{\rm D},i}H_i}\ge\sigma.
\]
The matching coefficient $q_i\ge w_{{\rm D},i}/(1+h_i)$ and
$H_i\le H_*$ complete the checks. Both independent assemblies verify
these inequalities piece by piece.

\begin{lemma}[Multiaffine allocation constants]\label{lem:joint-multiaffine}
Put $l=1/R$, $u=1/2$, and, for $0\le k\le n$, put
\[
 P_{n,k}=(1+l)^k(1+u)^{n-k},\quad
 E_{1,n,k}=kl+(n-k)u,\quad
 E_{2,n,k}=\binom k2l^2+k(n-k)lu+\binom{n-k}2u^2.
\]
Define
\[
 \eta_n=\min_{0\le k\le n}\frac{P_{n,k}-1}{E_{2,n,k}} (2\le n\le10),
 \qquad
 \zeta_n=\min_{0\le k\le n}\frac{P_{n,k}}{E_{1,n,k}} (1\le n\le9).
\]
Write $e_1(z)=\sum_i z_i$ and $e_2(z)=\sum_{i<j}z_iz_j$.
For $z_i\in[l,u]$,
$\prod_i(1+z_i)-1\ge\eta_n e_2(z)$ and
$\prod_i(1+z_i)\ge\zeta_n e_1(z)$ in the stated degree ranges.
In every degree the corresponding constants
$\eta_0=211/80$ and $\zeta_0=9/4$ are valid.
\end{lemma}

\begin{proof}
For a fixed trial constant each residual is multiaffine, so its minimum
on the box is attained at a vertex. The displayed formulas are precisely
its vertex tests, grouped by the number of lower coordinates.
For the uniform bounds enlarge the box to $[0,1/2]^n$.
At a vertex with $k$ nonzero coordinates the ratios are
$4((3/2)^k-1)/\binom k2$ and $2(3/2)^k/k$.
The first is minimized at $k=5$: its adjacent difference has the sign
of $(k-5)(3/2)^k+4$. The second is minimized at $k=2,3$,
as its adjacent difference has the sign of $k-2$.
The minima are $211/80$ and $9/4$.
\end{proof}

Use the normalized price assignments of
\cref{lem:four-routes-normalized}, repriced to the selected exponent $a$.
Let $x=A^{\rm arm}_F=A^{\rm arm}_E$ be the common raw-arm total of
\cref{def:closing-quantities}(3): the sum of the arm coefficients
before multiplying by their leading operator norms. Set $y=\sigma x$.
For this subsection abbreviate
$U=tU_0^{\prime\rm new}$ and $V=t{\Uinv}^{\rm new}$, where
$t=e^{2(a-a_0)570}$. These are the total first-stage price and the
portion assigned to invariant factors, respectively; $0\le V\le U$.
An outside factor means a factor on a block other than the selected
fold or matching blocks. Define
\begin{align*}
 F_\eta&=\eta_0e^U+\sum_{j=0}^8(\eta_{j+2}-\eta_0)V^j/j!,\\
 F_\zeta&=\zeta_0e^U+\sum_{j=0}^8(\zeta_{j+1}-\zeta_0)V^j/j!,\\
 B_\rho&=e^U+\sum_{j=1}^{10}(\rho^j-1)V^j/j!,\qquad \rho=1+1/R,\\
 w&=\frac{U}{1+\rho_H},\qquad
 \rho_H=\frac{3}{\lfloor\sqrt{Q_3}\rfloor(1-3/Q_3)}.
\end{align*}

\begin{proposition}[Four disjoint middle allocations]\label{prop:joint-middle-mean}
In the joint closing of \cref{cor:joint-closing}, the middle mean
decreases the closing number, in addition to its one-block mean
credits, by $\Delta_\alpha+\Delta_M+\Delta_O+\Delta_K$, where
\begin{align}
 \Delta_\alpha&=\frac{x^2}{16H_*^2}
                 \frac{e^w-1-w}{w^2},\label{eq:joint-alpha}\\
 \Delta_M&=\frac{x^2}{128}F_\eta,\qquad
 \Delta_O=\frac{x}{64}\left(y+\frac{y^3}{6}+\frac{y^5}{120}\right)F_\zeta,
 \qquad
 \Delta_K=\frac1{64}\left(\frac{y^2}{2}+\frac{y^4}{24}+\frac{y^6}{720}\right)B_\rho.
 \label{eq:joint-middle}
\end{align}
These are corrections to the combined comparison and mean bound;
they are not separate lower bounds obtained by substituting upper
prices in a mean energy.
\end{proposition}

\begin{proof}
After paying $R_{\rm err}e^{T_K+T_{QE}}$, the available positive product is
\[
 \bigotimes_B(\mu_B+E_{{\rm add},B}+G_B+H_B+K_{s,B}+Q_{E,B}).
\]
A tag specifies the summand chosen at a block in this tensor product.
Partition the retained terms first by their exact sets of $K_s$ tags
and then by their $Q_E$ tags. The four families are:
\begin{enumerate}[label=(\roman*)]
\item no $K_s$ or $Q_E$; at least two $\mathsf E$ tags, with reference
factors elsewhere;
\item no $K_s$; a nonempty $Q_E$ set and at least two positions carrying
$T$ or $Q_E$, with $\mu$ or $G$ at the remaining positions;
\item exactly $1,3,5$ $K_s$ tags; at least one further position carrying
$T$ or $Q_E$, with $\mu$ or $G$ at the remaining positions;
\item exactly $2,4,6$ $K_s$ tags, with arbitrary $\mu,T,Q_E,G$ outside.
\end{enumerate}
They are disjoint and correspond, in order, to
$\Delta_\alpha,\Delta_M,\Delta_O,\Delta_K$.
Each has support on at least two blocks, so none consumes a one-block
mean credit.

For the first family use
$\mathsf E_B\pm F_{{\rm fold},B}\succeq_{\vartheta_0}0$ on a selected pair.
The other $\mathsf E$ factors and reference factors are positive at
the same cut. The two companions dominate
$2\alpha_B(D_{F_1,B}+D_{F_2,B})/H_*$ after returning to the reserved
diagonal. Average the pair selection over the $\binom n2$ pairs of
an $n$-tag term. With raw coordinate $x_B=2\alpha_B$ this gives
coefficient $4/(H_*^2\binom n2)$ on the selected pair.
Consequently its exponential completion, after division by $64$, is
$x^2\sum_{j\ge0}w^j/(16H_*^2(j+2)!)$.

Here the $\mathsf E$ and $H$ assignments must be made together.
The physical bound $H_B\le\rho_H\mathsf E_B$ in positive mass allows
their total $U$ to be filled with an $\mathsf E$ assignment at least
$U/(1+\rho_H)$: first assign the genuine $H$ mass and put the remaining
slack into $\mathsf E$. The invariant and generic parts can be filled
with their original separate nonnegative slacks. This preserves $U$
and an invariant assignment $V_{\rm full}\ge V$. The completed
comparison cost is nonincreasing in $V_{\rm full}$, while the mean
corrections are nondecreasing in it, so replacing it by $V$ is safe.
No lower bound for the physical $\mathsf E$ mass is inferred from $U$.

For the second family factor out the $T$ coefficient at each matching
position and put $z_B=\alpha_B/\beta_B\in[1/R,1/2]$.
The nonempty-$Q_E$ sum dominates
$\prod T_B(\prod(1+z_B)-1)$.
The first inequality of \cref{lem:joint-multiaffine} allocates it
simultaneously over selected pairs; the factor $1/2$ in the unordered
pair completion gives $\Delta_M$.
For odd $K_s$ degree select one matching position instead and use the
second inequality. For even $K_s$ degree keep the whole outside
expansion and use $T_B+Q_{E,B}\succeq_\iota\rho T_B$ on a pure
invariant background. These give $\Delta_O$ and $\Delta_K$.
The uniform constants apply also to backgrounds containing $G$;
the finite improvements are used only on pure invariant backgrounds.
The multiaffine inequalities allocate a term among its eligible
positions; they never spend it separately in full on every position.

It remains to justify completion of the combined expression. For a
labelled diagonal energy $Z$, combine its comparison coefficient $c$
and retained mean coefficient $b$ as $(63c-b)Z$ first. The finite
checks below show $63c-b\ge0$, including the already retained variance
credits. Thus \cref{lem:record:record-expansion} applies to the residual
positive series. In particular, allowing repeated block labels enlarges
the sum over distinct blocks and gives the displayed exponential
series; this step is legitimate because the residual coefficients
are nonnegative. Its value is precisely
the former comparison bound minus the four displayed corrections.
This proves the proposition without asserting a lower bound for any
isolated upper-priced mean term.
\end{proof}

\paragraph{Residual coefficients.}
The sixteen classes and their invariant norm caps $\kappa_b$ are
those of \cref{cor:envelope-costs}. In the table below $F$ denotes a
selected fold arm and $E$ a selected matching arm. The outside degree
$j$ counts additional invariant factors. Check $j=0,\ldots,10$
separately, and use one further category for $j\ge11$.
The new mean coefficient is the sum of the applicable entries
\[
\begin{array}{c|c}
\text{selected class}&\text{mean coefficient}\\ \hline
F^2&4/H_*^2\\
E^2&\eta_{j+2}\ (j\le8),\quad\eta_0\ \text{otherwise}\\
F^nE,\ n=1,3,5&\sigma^n\zeta_{j+1}\ (j\le8),\quad\sigma^n\zeta_0\ \text{otherwise}\\
F^n,\ n=2,4,6&\sigma^n\rho^j\ (j\le10),\quad\sigma^n\ \text{in the tail}.
\end{array}
\]
The $F^2$ entries are added. Every selected block cap is at least $2$.
The restoration formulas therefore bound a selected-pair gain by $1/4$;
the same bound holds coefficientwise on each matching background.
Subtract at most $\kappa_b/4$ for $F^2$ and aligned $E^2$,
$\kappa_b/32$ for $F^4$, and the aligned $E^4$ cap divided by $32$.
All $16\cdot12=192$ residuals $63(\kappa_b-\text{variance})-
\text{mean}$ are positive; the minimum exceeds $23.66$.
The gates also hold before averaging over the record words. For the pure
fold classes $F^n$, $n=2,4,6$, the cap $\kappa_n$ is common to every word:
the chart reflection is independent of the word, and the extra reflection
acts trivially when there is no matching coordinate. In addition to the
class gates, check
\[
 63\bigl(1-\tfrac14\one_{n\in\{2,4\}}\bigr)\kappa_n
 -\frac4{H_*^2}\one_{n=2}-\sigma^n\rho^j>0
 \qquad(0\le j\le10),
\]
and the same inequality with $\rho^j$ replaced by $1$ for the tail.
These $36$ gates allow the full conservative variance credit on an
aligned quartic word, before its weight $1/8$ in the class is taken.
Each matching-pair class has a common cap; the pure matching quartics
have no new mean subtraction and retain their existing variance gates.
For the mixed words the chosen unframed caps are at least $1/4$,
and $63/4-\sigma^n\zeta_{j+1}>0$ for $n=1,3,5$, $0\le j\le8$.
The exact rational lists and their independent reconstruction are
part of \cref{cert:four-routes-closing}.

\begin{proposition}[Mean terms with a nonempty upper support]\label{prop:joint-upper-mean}
Let $S_u,A_u$ denote just the finite upper comparison assignments,
excluding prefix and tail padding. Set
\[
 C_E=e^{S_u}\cosh A_u-1,\qquad C_O=e^{S_u}\sinh A_u.
\]
The further joint correction is $\Delta_u-\mathcal R_u$, where
\begin{align*}
 \Delta_u={}&\frac1{64}\left\{
 C_E\bigl(\cosh(y)B_\rho+x\sinh(y)F_\zeta\bigr)
 +C_O\bigl(\sinh(y)B_\rho+x\cosh(y)F_\zeta\bigr)\right\},\\
 \mathcal R_u={}&\frac{R_{\rm err}e^{T_K+T_{QE}}}{64}
                   (e^{S_u+A_u}-1).
\end{align*}
\end{proposition}

\begin{proof}
On each finite upper block both matching brackets are $\mu_B+Q_B$,
with the same positive reserved-diagonal norm envelope $Q_B$.
The nonempty-$Q$ part of their tensor product is disjoint from all
middle-only terms used above. Fix its skew parity and the exact middle
$K_s$ set. If the parities agree, retain the full $T/Q_E/G$ background;
otherwise select one additional matching position. The same multiaffine
allocation gives the displayed even and odd exponential sums.
The new first-stage error is bounded by $\mathcal R_u$, obtained by
summing over the entire nonempty upper support, including discarded
branches. Every such term has positive upper degree, so it belongs to
none of the sixteen exceptional middle classes or their variance credits.
After splitting each full middle skew price into its raw arm and
positive remainder, the residual gates are
$63-\sigma^n\rho^j>0$ and $63-\sigma^n\zeta_{j+1}>0$, with the same
finite ranges and uniform tails. Since $\sigma<1$, $n=0$ is the worst
case, and the stated rational constants satisfy these gates.
Complete this positive residual series as above. The finite upper
assignments are used jointly in its positive cost and mean term;
they are not lower bounds for the mass of $Q_B$.
\end{proof}


\paragraph{Hypotheses verified in this section.}
\begin{itemize}
\item (CP6): for every atom the bracket of \cref{prop:record:matching}(i) holds.
 $M_{s,\Blk}-\Nbr_\Blk$ is a nonnegative combination of positive reserved-diagonal
 laws with uniform one-slot marginals, independent of the index: the norm
 envelopes $D^\epsilon_{\Blk,t}$ of (CP4) of the records other than the fold and
 matching records, the chosen matching diagonals $Q_{E,\Blk}$ and the chosen side
 diagonals $K_{s,\Blk}$. Its summed price is at most the $T$ of
 \cref{prop:record:matching}(ii), with the sides of \cref{prop:signed-sides}. The
 chosen diagonals are diagonal copies of components of the reference, positive
 records with admissible normalizations (\cref{def:record:frame} and the paragraph
 after it); so (CP8)(a) holds for them.
\item (CP7): $\frac12(e_{\rm L}+e_{\rm R})\ge(m-o(1))Y$ with $m$ of
 \eqref{eq:mean-credit}, by \cref{prop:record:mean-credit}.
\item Every bound used is an exact relation (orders at $\iota$ and $\vartheta_0$;
 the expansion and polarization identities; the contraction and density
 steps), an upper bound from \ref{it:record:E1}, or an exact relation minus such
 bounds (the restorations, the mean error), and the restorations are positive
 parts of continuous expressions; so the bound of (CP7) is robust, as (CP8)(b)
 requires (\cref{lem:record:rec}(i)). The gates of this section, $k^{\rm aux}_{s,i}\ge0$
 (positive in every piece), the protection $v_{\rm D}\ge2w_{\rm D}k_{\rm D}$, the
 matching gates, $n_Xe_0(X)<1$ and $c_{\ge2}(X)r_X<d_\tlaw(X)$, are inequalities
 between fixed numbers of the construction, which inflation does not change
 (\cref{rem:cp:robust-final}).
\end{itemize}
These are the rows ``(CP6)'' and ``(CP7)'' of \cref{tab:hypotheses}.

\paragraph{Consequences and what to read next.}
The arguments in this section use the following statements. The one-cut calculus and block-Gram tensorization
(\cref{lem:record:block-gram}); the sector lemma and signed mixtures
(\cref{lem:sector,rem:signed-mixture-sectors}); the admissible majorants
(\cref{lem:majorant}); the Paley triple calculus (\cref{lem:paley-triple});
polarization of all degrees
(\cref{lem:tr:polarization,rem:polarization-uses}); the energy toolkit
(\cref{lem:energy-toolkit}); the comparison-law interface and the robustness of
the final chain
(\cref{def:record:comparison-law,lem:record:record-expansion,lem:record:rec,rem:cp:robust-final,rem:cp:monotone});
repricing (\cref{lem:record:repricing,cor:first-stage-repricing}); and the
index-support and envelope-energy lemmas
(\cref{lem:record:common-mean,cor:index-support-frame,lem:envelope-energy}) with the
disjoint query endpoints of \cref{lem:slot-array}.

From the middle range: the frame, the first stage and the bracket with one fold
copy
(\cref{def:record:frame,lem:middle-polarization,lem:record:gram,lem:front12:reference},
\eqref{eq:protection}); the reference fold record, the actual features and the
signed law
(\cref{lem:old-f-law,lem:actual-features,def:record:mixture,lem:signed-components,prop:signed-bounds,lem:mirror},
\eqref{eq:signed-leading}, \eqref{eq:signed-constants}, \cref{cert:signed-local});
the prices of the matching, other middle and prefix records
(\cref{lem:record:exact-marginal-e,lem:middle-comparison-prices,lem:frontier:iidenergy});
the chamber geometry, restoration and coupons
(\cref{lem:record:fold-identification,lem:record:restoration,def:record:eligible,lem:record:paired-e-joint,lem:distinct-choice,lem:piece-caps,lem:pool-coupons})
with the piece data of \cref{subsec:variance:backgrounds,subsec:variance:final}; and
the block data
(\cref{def:record:exact-blocks,lem:record:local-domain,lem:record:block-charge}).

Finally: the index law (\cref{prop:oct24-index-cost}); the upper prices
(\cref{prop:record:upper}); the closing quantities
(\cref{def:closing-quantities}); and the certificates
\cref{cert:record:matching,cert:record:means,cert:closing}.
